% Generated draft: requires mathematical and source review.
\documentclass[11pt]{article}
\usepackage[margin=25mm]{geometry}
\usepackage{fontspec}
\setmainfont{Latin Modern Roman}
\usepackage{amsmath,amssymb,mathtools,unicode-math}
\setmathfont{Latin Modern Math}
\usepackage{enumitem,needspace,xurl,hyperref}
\setlength{\parindent}{0pt}
\setlength{\parskip}{5pt}
\setlength{\emergencystretch}{3em}
\newcommand{\N}{\mathbb{N}}
\newcommand{\Z}{\mathbb{Z}}
\newcommand{\Q}{\mathbb{Q}}
\newcommand{\R}{\mathbb{R}}
\newcommand{\C}{\mathbb{C}}
\newcommand{\F}{\mathbb{F}}
\newcommand{\A}{\mathbb{A}}
\newcommand{\PP}{\mathbb{P}}
\newcommand{\E}{\mathbb{E}}
\newcommand{\eps}{\varepsilon}
\newcommand{\dd}{\,\mathrm{d}}
\begin{document}
\section*{UnsolvedMath ID: 30003229}
\textbf{Source problem code:} catalogue entry
\subsection{Definitions and mathematical statement}
The supplied question asks whether the algebraic stringy Euler number strictly decreases under every Mori flip. Its precise domain is not specified in the record. The following gives a conventional complex projective interpretation; these domain assumptions require confirmation against the cited source.

For a smooth projective complex variety \(V\), let
\[
 b_{\mathrm{alg}}^{2i}(V)=\dim_{\Q}\operatorname{im}\bigl(\operatorname{cl}:\operatorname{CH}^{i}(V)_{\Q}\longrightarrow H^{2i}(V,\Q)\bigr),
 \qquad
 e_{\mathrm{alg}}(V)=\sum_{i=0}^{\dim V}b_{\mathrm{alg}}^{2i}(V),
\]
where \(\operatorname{CH}^{i}(V)_{\Q}\) is the group of codimension-\(i\) algebraic cycles modulo rational equivalence, with rational coefficients, and \(\operatorname{cl}\) is the cycle class map.

Suppose \(X\) is a normal projective complex variety with \(K_X\) \(\Q\)-Cartier and log terminal singularities. Choose a projective log resolution \(\pi:Y\to X\) with exceptional prime divisors \(D_1,\ldots,D_r\) having simple normal crossings, and write
\[
 K_Y=\pi^*K_X+\sum_{j=1}^{r}a_jD_j,\qquad a_j>-1.
\]
Put \(I=\{1,\ldots,r\}\), \(D_{\varnothing}=Y\), and
\[
 D_J=\bigcap_{j\in J}D_j,\qquad
 D_J^{\circ}=D_J\setminus\bigcup_{j\in I\setminus J}D_j
 \quad(J\subseteq I).
\]
For these strata use the additive extension
\[
 e_{\mathrm{alg}}(D_J^{\circ})
 =\sum_{J\subseteq K\subseteq I}(-1)^{|K|-|J|}e_{\mathrm{alg}}(D_K),
 \qquad e_{\mathrm{alg}}(\varnothing)=0.
\]
The algebraic stringy Euler number in this setting is
\[
 e_{\mathrm{alg}}^{\mathrm{str}}(X)
 =\sum_{J\subseteq I}e_{\mathrm{alg}}(D_J^{\circ})
   \prod_{j\in J}\frac{1}{a_j+1},
\]
with the empty product equal to \(1\).

In a conventional formulation, a Mori flip is a birational map \(g:X\dashrightarrow X^+\) fitting into small projective birational contractions \(f:X\to Z\) and \(f^+:X^+\to Z\), with \(f^+\circ g=f\) as rational maps, relative Picard numbers \(\rho(X/Z)=\rho(X^+/Z)=1\), \(-K_X\) relatively ample over \(Z\), and \(K_{X^+}\) relatively ample over \(Z\). Here small means that no prime divisor is contracted, and the relative Picard number is the dimension of the space of Cartier divisor classes modulo numerical equivalence over \(Z\).

\textbf{Question, conditional on this interpretation.} For every such Mori flip between normal projective complex \(\Q\)-Gorenstein varieties with log terminal singularities, does
\[
 e_{\mathrm{alg}}^{\mathrm{str}}(X)
 >e_{\mathrm{alg}}^{\mathrm{str}}(X^+)
\]
hold?

The supplied record does not explicitly impose the complex projective log terminal setting or specify whether \(\Q\)-factoriality is required. Consequently, this conditional formulation does not settle the intended scope of the source question. The record presents the inequality as an open question; its declarative original wording is not treated here as a proved theorem.
\subsection{Short English statement}
Does every Mori flip strictly decrease the algebraic stringy Euler number? A conventional complex projective log terminal interpretation is stated above, but the source record leaves the intended field, geometric hypotheses, and \(\Q\)-factoriality convention unspecified. The supplied open status requires independent verification.
\subsection{Sources}
\begin{itemize}
\item Supplied local UnsolvedMath registry record, permanent ID \(30003229\), title \emph{Stringy Euler Numbers under Mori Flips}, in the Oberwolfach Reports Open Problems collection. Exact source problem code: \texttt{catalogue entry}. The record gives proposed year 2016 and names no individual proposer.
\item Supplied source citation: \emph{Singularities (2017)}, Oberwolfach Reports, pp.\ 2625--2691, DOI \texttt{10.4171/owr/2016/46}. \url{https://doi.org/10.4171/owr/2016/46}. The background describes this report series as documenting mathematical workshops and their open-problem discussions.
\item Supplied literature reference: Batyrev--Gagliardi, \emph{On the algebraic stringy Euler number}. \url{https://arxiv.org/abs/1610.03842}.
\item Supplied literature reference: \emph{Oberwolfach Report 46}. \url{https://doi.org/10.4171/owr/2016/46}.
\item Supplied dated literature curation, checked on \texttt{2026-08-21}: the record labels the problem open and reports that strict decrease under divisorial Mori contractions or flips is proved for spherical and related equivariant classes, while no general proof or counterexample to the flip assertion was found. This assessment is supplied dated curation derived from the OWR record, not independently verified here; its time-sensitive status and cited claims require independent checking.
\item The registry labels the statement \texttt{corrected\_verified} and reports a manual semantic, typographic, and LaTeX review. This is supplied editorial metadata, not independent mathematical verification of this draft or of the conjectured inequality.
\end{itemize}
\end{document}
